\section{Contribution}

This article contributes a query-to-publication workflow scaffold for machine-learning research in the scientific process. The scaffold is not proposed as a replacement for field-specific reporting standards, statistical judgment, peer review, or domain expertise. Instead, it addresses the operational layer that often sits between those expectations and the final manuscript.

Existing reporting guidance commonly specifies what a completed manuscript should disclose. A workflow scaffold addresses a different question: how those disclosures are produced, reviewed, versioned, and linked to evidence during the research process. This distinction is especially important in machine-learning studies because weaknesses in the final paper often originate earlier in the workflow. These weaknesses include ambiguous research questions, model-first design, undocumented preprocessing, data leakage, weak baseline comparison, insufficient validation, and conclusions that exceed the strength of the evidence.

The proposed scaffold organizes machine-learning manuscript development into staged artifacts: research query, literature context, study protocol, data audit, feature plan, modeling plan, validation plan, result bundle, ethics review, reproducibility package, manuscript drafts, and submission package. Each artifact is designed to preserve a scientific decision that later supports the manuscript. The scaffold therefore treats writing as the synthesis of accumulated research artifacts rather than as a detached final step.

The contribution is procedural, reproducibility-oriented, and manuscript-aware. By linking question formation, protocol design, data documentation, feature engineering, model validation, ethical review, and publication preparation, the scaffold helps researchers align machine-learning practice with the logic of scientific inference. It also gives reviewers and collaborators a clearer trail from claim to evidence.

In this sense, the scaffold fills a practical gap between abstract reporting principles and project-specific code. It operationalizes reporting expectations as version-controlled workflow objects and manuscript modules, creating a repeatable structure for producing machine-learning research that is technically auditable and scientifically interpretable.
